
\section{\BookRecordIndex{T0397}在某处单射的曲线}\label{sec:2-5}
\zhindex{zaimouchudanshedequxian}{在某处单射的曲线}

设 $(\Sigma,j)$ 是紧黎曼曲面，$(M,J)$ 是殆复流形。如果存在紧黎曼曲面
$(\Sigma',j')$、\BookRecordIndex{T0003}$J$-全纯曲线 $u':\Sigma'\to M$ 以及\BookRecordIndex{T0287}全纯分歧覆盖映射
$\phi:\Sigma\to\Sigma'$，使得
\[
  u=u'\circ\phi,
  \qquad
  \deg(\phi)>1,
\]
则称 \BookRecordIndex{T0003}$J$-全纯曲线 $u:\Sigma\to M$ 是\emph{多重覆盖的}。
\zhindex{duozhongfugaiquxian}{多重覆盖曲线}
若 $u$ 不是多重覆盖的，则称 $u$ 是\emph{单纯的}。
\zhindex{danchunquxian}{单纯曲线}
我们将在下一章看到，对于一般的 $J$，给定同调类中的单纯 \BookRecordIndex{T0003}$J$-全纯曲线构成一个光滑
有限维流形。换言之，\BookRecordIndex{T0113}多重覆盖曲线是例外情形，它们可能成为 \BookRecordIndex{T0003}$J$-全纯曲线\BookRecordIndex{T0246}模空间中的
奇异点。这个结果的证明基于如下观察：每条单纯 \BookRecordIndex{T0003}$J$-全纯曲线都\emph{\BookRecordIndex{T0397}在某处单射}，
也就是说，存在某个 $z\in\Sigma$，使得
\[
  du(z)\ne0,
  \qquad
  u^{-1}(u(z))=\{z\}.
\]
具有这一性质的点 $z\in\Sigma$ 称为 $u$ 的一个\emph{\BookRecordIndex{T0094}单射点}。
\zhindex{danshedian}{单射点}
记
\[
  Z(u):=\bigl\{z\in\Sigma\bigm|du(z)=0
  \text{ 或 }\#u^{-1}(u(z))>1\bigr\}
\]
为\BookRecordIndex{T0094}单射点集的补集。下面的结果表明，每条单纯 \BookRecordIndex{T0003}$J$-全纯曲线的\BookRecordIndex{T0094}单射点集都是开且稠密的。

\begin{proposition}\label{prop:2-5-1}
设 $J$ 是 $M$ 上的 $C^2$ \BookRecordIndex{T0088}殆复结构，$\Sigma$ 是无边紧黎曼曲面，并且
$u:\Sigma\to M$ 是单纯 \BookRecordIndex{T0003}$J$-全纯曲线。则 $u$ \BookRecordIndex{T0397}在某处单射。此外，\BookRecordIndex{T0119}非单射点集 $Z(u)$
至多可数，并且只能在 $u$ 的\BookRecordIndex{T0233}临界点处聚积。
\end{proposition}

\begin{proof}
我们给出两个证明。它们都表明，任意曲线 $u:\Sigma\to M$ 都可以写成复合
\[
  u'\circ\phi:\Sigma\longrightarrow\Sigma'\longrightarrow M,
\]
其中 $u':\Sigma'\to M$ \BookRecordIndex{T0397}在某处单射。如果 $u$ 单纯，则 $\phi$ 的次数必为 $1$，从而
$u$ 也\BookRecordIndex{T0397}在某处单射。第一个证明从 $M$ 中的像曲线 $u(\Sigma)$ 构造定义域 $\Sigma'$；
第二个证明则借助 $\Sigma$ 上的一个等价关系构造 $\Sigma'$。

令 $X\subset\Sigma$ 表示 $u$ 的\BookRecordIndex{T0234}临界值的所有原像组成的集合，并令
\[
  X':=u(X)
\]
为\BookRecordIndex{T0234}临界值集。由引理~\ref{lem:2-4-1}，这两个集合都是有限的。令 $Q$ 为 $Y$ 中不同分支
相交之点的集合，其中
\[
  Y:=u(\Sigma)\setminus X'.
\]
这就是说，$x\in Q$ 当且仅当存在 $z_1,z_2\in\Sigma\setminus X$，使得
\[
  u(z_1)=u(z_2)=x,
  \qquad
  z_1\ne z_2,
  \qquad
  u(U_1)\ne u(U_2),
\]
其中最后一个不等式对 $z_j$ 的任意两个足够小的邻域 $U_j$ 都成立。由
引理~\ref{lem:2-4-3}，$Q$ 是 $Y$ 的离散子集（即它在 $Y$ 中没有聚点）。因此
\[
  S:=Y\setminus Q
\]
是 $M$ 中的嵌入子流形。以 $\iota:S\to M$ 表示这个嵌入。由于在 $Q$ 的每一点至多只有
有限多个 $Y$ 的分支相交，每个这样的点都给出 $S$ 的有限多个端，并且每个端都微分同胚于
穿孔圆盘 $\mathbb D\setminus\{\mathrm{pt}\}$。因此，可以向每个这样的端添加一个点，
并把 $\iota$ 光滑延拓为所得曲面 $S'$ 上的映射 $\iota'$。由于 $\iota'$ 是浸入，$S'$
上存在唯一的复结构，使 $\iota'$ 关于它是 $J$-全纯的。

流形 $S'$ 仍有与 $X'$ 中各点相对应的端。事实上，对每个 $x\in X'$，$\operatorname{Im}u$
可能有若干个经过 $x$ 的分支，而每个这样的分支都对应 $S'$ 的一个端。由此，每个端都是
穿孔圆盘的共形像，因而必具有穿孔圆盘的共形结构。所以，向 $S'$ 的每个端添加一点，
便可得到闭黎曼曲面 $\Sigma'$。又因为 $u$ 定义在整个 $\Sigma$ 上，映射 $\iota'$ 必可
延拓为 $J$-全纯映射 $u':\Sigma'\to M$。这个映射 $u'$ \BookRecordIndex{T0397}在某处单射，而且 $u$ \BookRecordIndex{T0213}可分解为
$u'\circ\phi$，其中 $\phi:\Sigma\to\Sigma'$ 是全纯映射。因此，$\phi$ 是分歧覆盖映射，
并且 $\phi$ 的次数为 $1$ 当且仅当 $u$ \BookRecordIndex{T0397}在某处单射。

下面给出一个由 $\Sigma$ 构造 $\Sigma'$ 的更详细论证。考虑集合
\[
  \Gamma_0\subset(\Sigma\setminus X)\times(\Sigma\setminus X),
\]
它由所有这样的偶对 $(z,\zeta)$ 组成：存在序列 $z_\nu\to z$ 和
$\zeta_\nu\to\zeta$，满足
$u(z_\nu)=u(\zeta_\nu)$、$z_\nu\ne z$ 且 $\zeta_\nu\ne\zeta$。即
\[
  \Gamma_0:=\bigl\{(z,\zeta)\bigm|
  z,\zeta\in\Sigma\setminus X,\qquad
  \exists z_\nu\to z\ \exists\zeta_\nu\to\zeta:\qquad
  u(z_\nu)=u(\zeta_\nu),\ z_\nu\ne z,\ \zeta_\nu\ne\zeta\bigr\}.
\]
换言之，$\Gamma_0$ 是 $u$ 的重合点的聚点集；孤立的自交点被排除在外。我们来证明
$\Gamma_0$ 具有下列性质。
\begin{enumerate}[label=\textup{(\roman*)}]
  \item $\Gamma_0$ 是 $\Sigma\setminus X$ 上的等价关系。
  \item 到第一因子的投影 $\pi:\Gamma_0\to\Sigma\setminus X$ 是局部微分同胚，且
  $\pi$ 的每个局部逆都是全纯的。
  \item 覆盖映射 $\pi$ 是恰当的。
  \item 对每个 $z\in\Sigma\setminus X$，$\pi^{-1}(z)$ 都是有限集，并且数目
  \[
    m:=\#\bigl\{\zeta\in\Sigma\setminus X\bigm|(z,\zeta)\in\Gamma_0\bigr\}
  \]
  与 $z\in\Sigma\setminus X$ 无关。
\end{enumerate}

性质 (i) 来自引理~\ref{lem:2-4-3}。事实上，若
$(z_0,z_1)\in\Gamma_0$ 且 $(z_1,z_2)\in\Gamma_0$，则由
引理~\ref{lem:2-4-3}，存在 $z_i$ 的开邻域 $U_i\subset\Sigma$（$i=0,1,2$）和
全纯映射 $\phi_{10}:U_0\to U_1$、$\phi_{21}:U_1\to U_2$，使得
\[
  \phi_{10}(z_0)=z_1,
  \qquad
  \phi_{21}(z_1)=z_2,
  \qquad
  \left.u\right|_{U_0}=u\circ\phi_{10},
  \qquad
  \left.u\right|_{U_1}=u\circ\phi_{21}.
\]
映射 $\phi_{10}$ 和 $\phi_{21}$ 显然都不是常值映射。由于
\[
  \left.u\right|_{U_0}=u\circ\phi_{21}\circ\phi_{10},
  \qquad
  \phi_{21}\circ\phi_{10}(z_0)=z_2,
\]
可知 $(z_0,z_2)\in\Gamma_0$。这就证明了 (i)。性质 (ii) 来自同一个论证，以及如下
事实：因为 $z_0\notin X$，全纯映射 $\phi_{10}$ 满足 $d\phi_{10}(z_0)\ne0$。
为证明 (iii)，设 $z_\nu\in\Sigma\setminus X$ 是收敛到
$z\in\Sigma\setminus X$ 的序列，且 $(z_\nu,\zeta_\nu)\in\Gamma_0$。我们必须证明
$\zeta_\nu$ 有一个收敛到 $\Sigma\setminus X$ 中某点的子列。取子列后，不妨设
$\zeta_\nu$ 收敛到 $\zeta\in\Sigma$。此时 $u(z)=u(\zeta)$，因而
\[
  \zeta\in u^{-1}\bigl(u(\Sigma\setminus X)\bigr)=\Sigma\setminus X.
\]
这就说明 $\pi$ 是恰当的。性质 (iv) 来自 (ii)、(iii) 以及
$\Sigma\setminus X$ 连通这一事实。

现在的策略是证明：若 $m\ge2$，则 $u$ 不是单纯的；若 $m=1$，则 $u$ 满足命题的要求。
先设 $m=1$。我们必须证明，集合 $Z(u)$ 只能在 $u$ 的\BookRecordIndex{T0233}临界点集上聚积。反设存在序列
$z_\nu\in Z(u)$ 收敛到某个点 $z_0\in\Sigma$，且 $du(z_0)\ne0$。选择 $z_0$ 的开邻域
$U_0\subset\Sigma$，使 $u$ 在 $U_0$ 上的限制是嵌入。不妨设对每个 $\nu$ 都有
$z_\nu\in U_0$。于是对每个 $\nu$ 都有 $du(z_\nu)\ne0$，从而存在序列
$\zeta_\nu\in\Sigma\setminus\{z_\nu\}$，满足 $u(z_\nu)=u(\zeta_\nu)$。因为
$\left.u\right|_{U_0}$ 是嵌入，所以 $\zeta_\nu\notin U_0$。不妨设 $\zeta_\nu$ 收敛到
某个 $\zeta_0\in\Sigma\setminus U_0$。于是由引理~\ref{lem:2-4-3}，存在
$\zeta_0$ 的开邻域 $V_0\subset\Sigma$ 和全纯映射 $\phi:V_0\to U_0$，使得
$\phi(\zeta_0)=z_0$ 且
\[
  \left.u\right|_{V_0}=u\circ\phi.
\]
必要时缩小 $V_0$，可以假设
$V_0\setminus\{\zeta_0\}\subset\Sigma\setminus X$。于是，对每个
$\zeta\in V_0\setminus\{\zeta_0\}$ 都有 $(\zeta,\phi(\zeta))\in\Gamma_0$，从而
$m\ge2$，这与 $m=1$ 的假设矛盾。这证明了当 $m=1$ 时，集合 $Z(u)$ 只能在 $u$ 的
\BookRecordIndex{T0233}临界点集上聚积。

现在设 $m\ge2$，并把 $\Gamma_0$ 扩充到它的闭包
\[
  \Gamma:=\operatorname{cl}(\Gamma_0)\subset\Sigma\times\Sigma.
\]
这个集合是 $\Sigma$ 上的等价关系；我们记
\[
  z\sim\zeta
  \qquad\Longleftrightarrow\qquad
  (z,\zeta)\in\Gamma.
\]
集合 $X$ 在这个等价关系下保持不变。（若 $z\in X$ 且 $z\sim\zeta$，则
$\zeta\in X$。）因此，每个点 $z\in\Sigma$ 都带有一个自然的重数 $m(z)\ge1$，其定义
如下。若 $z\in\Sigma\setminus X$，定义 $m(z)=1$。若 $z\in X$，而
$w\in\Sigma\setminus X$ 充分接近 $z$，则 $w$ 的等价类中的全部 $m$ 个点都接近 $X$。
定义 $m(z)$ 为与 $w$ 等价且接近 $z$ 的点数。由连续性，这个数与 $w$ 的选取无关。按此
定义，对每个 $z\in\Sigma$ 都有
\[
  \sum_{\zeta\sim z}m(\zeta)=m.
\]

把拓扑空间 $\Sigma'$ 定义为商空间
\[
  \Sigma':=\Sigma/{\sim}.
\]
还需要在 $\Sigma'$ 上找到一个转移映射全纯的坐标图册，使自然投影
$\phi:\Sigma\to\Sigma'$ 是全纯的。若 $z_0\in\Sigma\setminus X$，则 $[z_0]$ 附近所需
坐标图，只须把 $\Sigma$ 上的一个全纯坐标图限制到 $z_0$ 的一个足够小邻域。现在设
$z_0\in X$，并记 $m_0:=m(z_0)$。于是，对于 $z_0$ 附近的每个 $z$，在 $z_0$ 附近恰有
$m_0$ 个互不相同且与 $z$ 等价的点（包括 $z$ 本身）。$z$ 的等价类中其余各点都靠近
$X$ 的其他点。选择 $z_0$ 的邻域 $U_0\subset\Sigma$，使它在等价关系的这个局部部分下
保持不变。那么，对每个 $z\in U_0\setminus\{z_0\}$，$U_0$ 与 $z$ 的等价类之交恰含
$m_0$ 个点。必要时缩小 $U_0$，可以假设存在全纯坐标图 $w:U_0\to\mathbb C$，满足
\[
  w(z_0)=0.
\]
于是，$\Sigma'$ 上所需的坐标图是映射 $w':U_0/{\sim}=:U_0'\to\mathbb C$，其中
\[
  w'([z]):=\prod_{\substack{\zeta\in U_0\\ \zeta\sim z}}w(\zeta).
\]
更准确地说，由标准的提升论证和单复变量全纯映射的\BookRecordIndex{T0214}可去奇点定理，存在 $m_0$ 个全纯
微分同胚
\[
  \psi_j:U_0\to U_0,
  \qquad j=1,\ldots,m_0,
\]
使得 $\psi_1=\mathrm{id}$，每个 $\psi_j$ 都满足 $\psi_j(z_0)=z_0$，并且对
$z,\zeta\in U_0$ 有
\[
  z\sim\zeta
  \qquad\Longleftrightarrow\qquad
  \zeta\in\{\psi_1(z),\ldots,\psi_{m_0}(z)\}.
\]
采用这个记号，映射 $w':U_0'\to\mathbb C$ 可写成
\[
  w'([z])=\prod_{j=1}^{m_0}w\circ\psi_j(z).
\]
因此，从 $U_0$ 到 $U_0'$ 的投影与 $w':U_0'\to\mathbb C$ 的复合是全纯映射，并且在
$z_0$ 有 $m_0$ 阶零点。所以，$w'$ 像中的一个非零元素在上述复合映射下恰有 $m_0$ 个
原像，它们正是 $U_0$ 上局部等价关系的一个等价类中的元素。因此
$w':U_0'\to\mathbb C$ 是到其像上的同胚。转移映射的全纯性来自等价关系 $\Gamma_0$
的上述性质 (ii)。

由构造，$\phi:\Sigma\to\Sigma'$ 是次数为 $m$ 的全纯映射，且
\begin{equation}\label{eq:2-5-1}
  z\sim\zeta
  \qquad\Longleftrightarrow\qquad
  \phi(z)=\phi(\zeta).
\end{equation}
因此存在唯一映射 $u':\Sigma'\to M$，使得 $u=u'\circ\phi$。映射 $u'$ 在 $\Sigma'$ 上
显然连续，并在有限集 $\phi(X)$ 的补集上是 $J$-全纯的。它在 $\phi(X)$ 中每一点附近的
能量都是有限的；于是，由 \BookRecordIndex{T0003}$J$-全纯曲线的\BookRecordIndex{T0214}可去奇点定理（见下面的定理 4.1.2），$u'$ 是
光滑的。由于 $\deg(\phi)=m\ge2$，可知 $u$ 不是单纯的。
\end{proof}

\begin{remark}\label{rem:2-5-2}
假设在命题~\ref{prop:2-5-1} 的证明中 $\Sigma=\mathbb{CP}^1$。那么 $\Sigma'$ 同构于
$\mathbb{CP}^1$；在这个同构下，全纯函数 $\phi:\Sigma\to\Sigma'$ 对应于一个有理函数，
仍记为 $\phi:\mathbb{CP}^1\to\mathbb{CP}^1$。下面给出从全纯等价关系 $\Gamma$ 构造
有理函数的一个直接论证，并使式~\eqref{eq:2-5-1} 成立。把 $\mathbb{CP}^1$ 与黎曼球面
$\mathbb C\cup\{\infty\}$ 等同，并使 $0\not\sim\infty$，记
\[
  P_\infty:=\{z\in\mathbb C\cup\{\infty\}\mid z\sim\infty\},
  \qquad
  P_0:=\{z\in\mathbb C\cup\{\infty\}\mid z\sim0\}.
\]
于是 $\infty\in P_\infty$、$0\in P_0$，而 $0\notin P_\infty$、
$\infty\notin P_0$。定义 $\phi:\mathbb C\setminus P_\infty\to\mathbb C$ 为
\begin{equation}\label{eq:2-5-2}
  \phi(z):=\prod_{\zeta\sim z}\zeta^{m(\zeta)},
  \qquad z\in\mathbb C\setminus P_\infty.
\end{equation}
那么 $\phi$ 是全纯的，并可连续延拓为从 $\mathbb C\cup\{\infty\}$ 到自身的函数。
（请证明这一点！）因此 $\phi$ 是有理函数。它在每个 $z\in P_0$ 处都有 $m(z)$ 阶零点，
并且没有其他零点，所以 $\phi$ 的次数是 $m$。由定义，
$z\sim\zeta$ 蕴含 $\phi(z)=\phi(\zeta)$。反向蕴含则来自 $\phi$ 的次数为 $m$ 这一事实。
因此 $\phi$ 满足式~\eqref{eq:2-5-1}。
\end{remark}

\begin{exercise}\label{ex:2-5-3}
设 $u:\mathbb C\cup\{\infty\}\to\mathbb C\cup\{\infty\}$ 是非常值有理函数，并且
$u(0)\ne u(\infty)$。在黎曼球面上定义等价关系
\[
  z\sim\zeta
  \qquad\Longleftrightarrow\qquad
  u(z)=u(\zeta).
\]
定义函数 $m:\mathbb C\cup\{\infty\}\to\mathbb N$ 如下：对 $u$ 的每个极点 $z$，令
$m(z)$ 为 $z$ 作为极点的阶数；对其他每个点 $z$，令 $m(z)$ 为 $z$ 作为 $u-u(z)$ 的
零点的阶数。令 $\phi$ 由式~\eqref{eq:2-5-2} 给出。证明存在一个 \BookRecordIndex{T0040}Möbius 变换 $T$，使得
$u=T\circ\phi$。
\end{exercise}

\begin{corollary}\label{cor:2-5-4}
设 $J$ 是 $M$ 上的 $C^2$ \BookRecordIndex{T0088}殆复结构，$\Sigma_0,\Sigma_1$ 是无边紧连通黎曼曲面，并且
$u_j:\Sigma_j\to M$ 是单纯 \BookRecordIndex{T0003}$J$-全纯曲线，满足
\[
  u_0(\Sigma_0)=u_1(\Sigma_1).
\]
则存在全纯微分同胚 $\phi:\Sigma_1\to\Sigma_0$，使得
\[
  u_1=u_0\circ\phi.
\]
\end{corollary}

\begin{proof}
对 $i=0,1$，令 $Z_i\subset\Sigma_i$ 为 $u_i$ 的\BookRecordIndex{T0119}非单射点集。于是 $Z_i$ 由 $u_i$ 的
\BookRecordIndex{T0233}临界点以及所有满足下述条件的点 $z\in\Sigma_i$ 组成：$u_i^{-1}(u_i(z))$ 还含有 $z$
以外的点。特别地，$u_i$ 的\BookRecordIndex{T0234}临界值的每个原像，只要其本身不是\BookRecordIndex{T0233}临界点，就满足第二个条件，
因而属于 $Z_i$。所以
\[
  Z_0=u_0^{-1}(u_0(Z_0)),
  \qquad
  Z_1=u_1^{-1}(u_1(Z_1)).
\]
由命题~\ref{prop:2-5-1} 和引理~\ref{lem:2-4-1}，$Z_i$ 至多可数，并且只能在由
$u_i$ 的\BookRecordIndex{T0233}临界点组成的有限子集处聚积。由于 $u_0(\Sigma_0)=u_1(\Sigma_1)$，存在唯一的
双射
\[
  \phi:
  \Sigma_1\setminus u_1^{-1}\bigl(u_0(Z_0)\cup u_1(Z_1)\bigr)
  \longrightarrow
  \Sigma_0\setminus u_0^{-1}\bigl(u_0(Z_0)\cup u_1(Z_1)\bigr),
\]
使得 $u_1=u_0\circ\phi$。对 $\phi$ 定义域中的每个 $z_1$，
$du_0(\phi(z_1))$ 和 $du_1(z_1)$ 都非零且复线性；由链式法则，$\phi$ 是全纯微分同胚。
此外，由引理~\ref{lem:2-4-1}，集合
$u_1^{-1}\bigl(u_0(Z_0)\cup u_1(Z_1)\bigr)$ 至多可数，并且只能在有限多个点处聚积。
在合适的局部坐标中表示时，$\phi$ 是有界的；因此，单复变量全纯函数的\BookRecordIndex{T0214}可去奇点定理表明，
$\phi$ 可延拓为从 $\Sigma_1$ 到 $\Sigma_0$ 的全纯微分同胚。这就证明了推论。
\end{proof}

\begin{corollary}\label{cor:2-5-5}
设 $J$ 是 $M$ 上的 $C^2$ \BookRecordIndex{T0088}殆复结构，$\Sigma_0,\ldots,\Sigma_N$ 是无边紧连通黎曼曲面，
并且对 $j=0,\ldots,N$，$u_j:\Sigma_j\to M$ 是 \BookRecordIndex{T0003}$J$-全纯曲线。假设 $u_0$ 单纯且
\[
  u_0(\Sigma_0)\ne u_j(\Sigma_j)
  \qquad\text{对 }j>0.
\]
那么，对每个 $z_0\in\Sigma_0$ 以及 $z_0$ 的每个开邻域 $U_0\subset\Sigma_0$，都存在
一个以 $z_0$ 为中心的\BookRecordIndex{T0170}环域 $A_0\subset U_0$，使得 $u_0:A_0\to M$ 是嵌入，并且
\[
  u_0^{-1}(u_0(A_0))=A_0,
  \qquad
  u_0(A_0)\cap u_j(\Sigma_j)=\varnothing
  \qquad\text{对 }j>0.
\]
\end{corollary}

\begin{proof}
对 $j>0$，令
\[
  Z_j:=u_0^{-1}(u_j(\Sigma_j)),
\]
并令 $Z_0\subset\Sigma_0$ 表示 $u_0$ 的\BookRecordIndex{T0119}非单射点集。由命题~\ref{prop:2-5-1}，集合
$Z_0$ 只能在 $u_0$ 的\BookRecordIndex{T0233}临界点处聚积；由命题~\ref{prop:2-4-4}，集合
$Z_1,\ldots,Z_N$ 也具有同样
的性质。在 $U_0$ 中选取一个以 $z_0$ 为中心、包含于
$Z_0\cup\cdots\cup Z_N$ 的补集的\BookRecordIndex{T0170}环域 $A_0$。那么 $A_0$ 满足本推论的要求。
\end{proof}

\begin{remark}\label{rem:2-5-6}
在许多情形中，考虑以\BookRecordIndex{T0219}拉格朗日子流形 $L$ 为边界条件的 $J$-全纯圆盘
\[
  u:(\mathbb D,\partial\mathbb D)\longrightarrow(M,L)
\]
很有用。对这类对象，需要一个精细得多的“多重覆盖”概念，因为圆盘可能只部分地绕着自身
覆盖：可以设想 $S^2$ 中一个以赤道为边界、绕球面一周半的圆盘。这个问题最终由
Lazzarini \srcite{234, 235} 澄清。另见 Kwon--Oh \srcite{220} 和
Zehmisch \srcite{425}。
\end{remark}
